\section*{Chapitre \ref{ch:ll:vector-instructions-and-data-parallel-code} --- Instructions vectorielles et code à parallélisme de données}

\begin{solution}{exo:ll:vector-instructions-and-data-parallel-code:1}
4 doubles, 8 entiers de 32 bits, 32 octets. $1\,500 = 46 \times 32 + 28$~: 46 blocs complets et un dernier bloc qui chevauche le précédent, soit 47
comparaisons.
\end{solution}

\begin{solution}{exo:ll:vector-instructions-and-data-parallel-code:2}
$\mathtt{0x00041010}$ a les bits 4, 12 et 18 à 1~: délimiteurs aux octets 4, 12 et 18 du bloc.
\end{solution}

\begin{solution}{exo:ll:vector-instructions-and-data-parallel-code:3}
Après la première étape, les voies de 16 bits contiennent 20, 26, 9 et 25~; après la deuxième, les voies de 32 bits contiennent 2\,026 et 925~;
la troisième étape donne $2\,026 \times 10\,000 + 925 = 20\,260\,925$.
\end{solution}

\begin{solution}{exo:ll:vector-instructions-and-data-parallel-code:4}
Octets 0--31 puis 32--63, puis le dernier chargement couvre 38--69. Après deux blocs complets, la position est 64~; le dernier masque est donc décalé
vers la droite de $64 - 38 = 26$ bits~: ses bits restants correspondent aux octets 64--69. Un délimiteur à l'octet 60 est trouvé dans le deuxième bloc,
avant le dernier chargement~; le décalage ne fait que retirer des octets déjà examinés.
\end{solution}

\begin{solution}{exo:ll:vector-instructions-and-data-parallel-code:5}
Aliasing possible~: la sortie pourrait chevaucher une entrée, donc la boucle est versionnée par une vérification à l'exécution~; \texttt{\_\_restrict}
promet l'absence de recouvrement et supprime cette vérification. Pas de type vectoriel pour les champs de l'enregistrement~: les champs sont séparés de 80 octets,
on les stocke donc en colonnes. Flot de contrôle dans la boucle~: la recherche peut se terminer tôt~; on l'écrit avec des intrinsèques, ou on parcourt
un bloc entier avec un masque comme le fait la construction.
\end{solution}

\begin{solution}{exo:ll:vector-instructions-and-data-parallel-code:6}
Environ \qty{52}{\nano\second} sur la mesure retenue, soit près de 79 octets par nanoseconde. Dès que le tampon dépasse le cache de premier niveau,
c'est le débit auquel les niveaux suivants fournissent les lignes qui limite le calcul, et non les comparaisons.
\end{solution}

\begin{solution}{exo:ll:vector-instructions-and-data-parallel-code:7} \sloppy
Recopier \texttt{positions\_any\_avx2} avec des registres de 128 bits (\texttt{\_mm\_cmpeq\_epi8} et
\texttt{\_mm\_movemask\_epi8}) et des blocs de 16 octets (le dernier commence à $n - 16$, décalé de $i - (n - 16)$), revenir à la boucle scalaire
en dessous de 16 octets, faire en sorte que la répartition renvoie cette version à \texttt{Level::sse2}, et ajouter une comparaison avec
\texttt{positions\_any\_scalar} dans la boucle sur les longueurs du test.
\end{solution}

\begin{solution}{exo:ll:vector-instructions-and-data-parallel-code:8}
\texttt{-march=native} compile pour le processeur de la machine de construction, AVX-512 compris, partout où le compilateur choisit de l'utiliser~;
les exécuteurs d'intégration continue et les nouveaux serveurs le possédaient aussi, par hasard. Le premier hôte de production plus ancien exécute
une instruction qu'il n'a pas, et le processus se termine par un signal d'instruction illégale. Il faut compiler pour le processeur le plus ancien
du parc et répartir le code plus large à l'exécution.
\end{solution}

\begin{solution}{pb:ll:vector-instructions-and-data-parallel-code:1}
\begin{enumerate}
\item Environ \qty{0.22}{\nano\second} par octet pour la boucle scalaire, \qty{0.024}{\nano\second} pour SSE2 et \qty{0.013}{\nano\second}
pour AVX2.
\item Sur l'exécution retenue, $a \approx \qty{0.3}{\nano\second}$ et $b \approx \qty{0.0125}{\nano\second}$ par octet.
\item $L^* = a/(b_{\text{scalaire}} - b) \approx 2$ octets pour AVX2 et environ 8 pour SSE2.
\item La boucle scalaire paie aussi un coût par appel (l'appel, le branchement de sortie)~; sa droite se situe donc au-dessus de celle qui passe par l'origine,
et le point de croisement réel est plus tôt.
\item En ramenant la répartition mesurée à 200 octets~: environ \qty{120}{\nano\second} octet par octet, \qty{14}{\nano\second} avec AVX2.
\item $200\,000 \times \qty{120}{\nano\second} \approx \qty{24}{\milli\second}$ par seconde (2.4~\% d'un cœur) contre
\qty{2.8}{\milli\second}.
\item $4 \times (42.6 - 6.6) \approx \qty{144}{\nano\second}$ par rapport.
\item L'analyse syntaxique~: \qty{144}{\nano\second} contre environ \qty{106}{\nano\second} pour le découpage.
\item Il n'y a rien à chercher~: les champs sont lus à des décalages connus.
\item La répartition appelle la version SSE2, que tout processeur x86-64 exécute~; rien ne échoue.
\item Une lecture au-delà de la fin du message peut franchir dans une page non mappée et provoquer une faute, et elle lit des octets qui ne sont pas ceux
du message~; le bloc qui chevauche ne lit que le message.
\item Deux fois plus d'octets par comparaison, donc à peu près moitié moins de coût par octet sur les longs tampons~; rien pour les messages courts, où le
coût fixe domine.
\item \textbf{Résultat nommé.} La recherche AVX2 coûte environ \qty{0.3}{\nano\second} par appel plus \qty{0.0125}{\nano\second} par
octet, contre \qty{0.22}{\nano\second} par octet pour la boucle scalaire~: elle amortit sa mise en place à partir de 2 octets (SSE2 à partir de
8 environ), donc aucun message n'est trop court pour elle.
\item Des champs isolés peuvent l'être (\texttt{54=1} fait cinq octets avec son délimiteur), mais le balayage porte sur l'ensemble du message, jamais
sur un seul champ.
\item La même mesure sur le processeur du serveur, avec la distribution réelle des longueurs de message et des caches froids.
\item Ils définissent la bonne réponse, à laquelle les tests comparent chaque version vectorielle, et ils servent de repli.
\item Appeler chaque version directement dans les tests, comme le fait la construction, et ajouter un moyen de forcer le niveau (une variable d'environnement
lue par \texttt{detect}) afin que le programme entier puisse s'exécuter sur le repli sur une machine qui possède AVX2.
\item Sommes de contrôle, recopie, analyse JSON et FIX, recherches sur les niveaux de prix, et calculs groupés sur les positions.
\item Garder les champs que lisent les contrôles (prix, quantités, limites) en colonnes.
\item Lorsque la même opération s'applique à de nombreux éléments contigus.
\end{enumerate}
\end{solution}

\begin{solution}{iq:ll:vector-instructions-and-data-parallel-code:1}
Une instruction appliquée à plusieurs éléments dans un registre large. Une boucle qui additionne deux tableaux de doubles élément par élément en tire profit~;
parcourir une liste chaînée, ou une boucle dont chaque itération dépend du résultat de la précédente, n'en tire pas profit.
\iqlookfor{itérations indépendantes sur des données contiguës.}
\end{solution}

\begin{solution}{iq:ll:vector-instructions-and-data-parallel-code:2}
Comparer 16 ou 32 octets à la fois, l'octet recherché étant diffusé dans un registre vectoriel, prendre le masque des octets égaux et compter ses zéros de fin~;
terminer par un dernier bloc qui chevauche plutôt que de lire au-delà de la fin~; choisir la largeur à l'exécution.
Vérifier contre une boucle scalaire pour chaque longueur.
\iqlookfor{comparer, masquer, compter et traiter la queue.}
\end{solution}

\begin{solution}{iq:ll:vector-instructions-and-data-parallel-code:3}
Aliasing possible, sorties anticipées, données non contiguës ou à pas, appels non insérés, réductions en virgule flottante dont le langage fixe l'ordre,
et modèle de coût qui juge la vectorisation non rentable. On interroge le compilateur~: \texttt{-fopt-info-vec-missed} dans GCC donne la raison boucle par boucle~;
on lit ensuite l'assembleur.
\iqlookfor{obstacles concrets et option de rapport.}
\end{solution}

\begin{solution}{iq:ll:vector-instructions-and-data-parallel-code:4}
Compiler pour le processeur le plus ancien, placer le code AVX2 dans des fonctions compilées pour AVX2 par un attribut (ou par \omterm{def:ll:cpp-for-latency-iii-the-compiler:fmv}{multiversionnage de fonctions}),
et choisir à l'exécution selon les fonctions que rapporte le processeur~; tester les deux chemins.
\iqlookfor{répartition à l'exécution, et non drapeau global.}
\end{solution}

\begin{solution}{iq:ll:vector-instructions-and-data-parallel-code:5}
\omterm{def:ll:vector-instructions-and-data-parallel-code:soa}{Structure de tableaux} pour la revalorisation~: chaque grecque des positions consécutives est contiguë, donc elle remplit les registres vectoriels et
ne charge que les champs lus. Les enregistrements restent préférables pour le code qui traite une position à la fois~; un hybride de petits blocs
sert les deux usages.
\iqlookfor{le motif d'accès décide de la disposition.}
\end{solution}

\begin{solution}{iq:ll:vector-instructions-and-data-parallel-code:6}
Découper au point décimal, analyser chaque partie huit chiffres à la fois dans un registre de 64 bits (soustraire le caractère zéro,
puis combiner les paires, les quadruplets et les octets par trois étapes de multiplication et masquage), et mettre à l'échelle la fraction avec les décimales manquantes.
Valider que chaque octet est un chiffre (un test de bits sur le mot entier), le nombre de décimales, l'intervalle et le signe.
\iqlookfor{les étapes de multiplication et masquage et la validation.}
\end{solution}
